\subsection{问题二：历史预测、风险计划与实时执行}
\subsubsection{历史负荷与光伏预测}
设$P^L_{d,t},P^V_{d,t}$为功率观测。对具有至少一周历史的日期，负荷的同星期预测为
\begin{equation}
 \widehat P^L_{d,t}=\frac{\sum_{k=1}^{K_d}0.55^{k-1}P^L_{d-7k,t}}
 {\sum_{k=1}^{K_d}0.55^{k-1}},\qquad K_d=\min\{4,\lfloor d/7\rfloor\}.
\end{equation}
日期索引从1月1日的$d=0$开始。历史不足7天时，改用此前最多7天同刻值，以$0.8^{k-1}$加权。光伏采用较短窗口：
\begin{equation}
 \widehat P^{V,H}_{d,t}=\frac{\sum_{k=1}^{\min(7,d)}0.75^{k-1}P^V_{d-k,t}}
 {\sum_{k=1}^{\min(7,d)}0.75^{k-1}}.
\end{equation}
无历史的首日负荷先验为4500\,kW、光伏为0，仅用于1月初始运行。两类预测均只读取日期$d$之前的观测，不把附件一包含全年信息的典型曲线作为先验。净负荷电量点预测为
\begin{equation}
 \widehat N_{d,t}=\Delta(\widehat P^L_{d,t}-\widehat P^{V,H}_{d,t}),\qquad\Delta=1/6\ \mathrm{h}.
\end{equation}

\subsubsection{经验风险分位与零时计划}
直接按点预测购电会低估紧急购电的非对称代价。对预测发布时刻$r$，保存其对应预测器的历史残差
\begin{equation}
 R^{(r)}_{j,t}=N_{j,t}-\widehat N^{(r)}_{j,t}.
\end{equation}
仅使用$\max(7,d-28)\le j<d$的日期，将同刻及前后各3个十分钟段的残差合并；在当前发布时刻至日末的剩余预测区间两端，按端点复制补齐邻域。由所得集合$\mathcal R^{(r)}_{d,t}$形成风险调整量
\begin{equation}
 \widetilde N^{(r)}_{d,t}=\widehat N^{(r)}_{d,t}
   +Q_{\alpha}\!\left(\mathcal R^{(r)}_{d,t}\right).\label{eq:risk}
\end{equation}
经验分位对排序样本采用线性插值，即NumPy的\texttt{quantile}默认规则。没有足够历史时不添加分位修正。不同光伏融合权重或信息子集各自重算历史预测与残差，避免混用预测器的误差分布。

零时读取上一日真实末库存，在当日144段上将式\eqref{eq:balance}的净负荷换成$\widetilde N$、令$e_t=0$，最小化预测常规购电费，增加$S_{144}\ge6000$的规划约束。由此生成固定原计划$g_t$。分位参数越大，通常越倾向于提前购入电量，但可能产生更多付费余电；费用随参数未必单调，且电池把各时段连接起来，不能把单段分位直接解释为全系统的供电可靠性保证。

已有两阶段微网调度研究将日前储能SOC下限传入日内决策，并联合考虑应急供能成本\cite{hu2025emergency}，说明库存管理与应急支出需要协调。本文的6000\,kWh是自行设定的日末规划恢复目标，并非从该文推得的最优下限，也不涉及其应急电源车决策。

候选分位为$\{0.50,0.65,0.70,0.80,0.90,0.95\}$，按1月15日至31日的实际现金支出选择；候选从1月1日6000\,kWh起以自身固定参数回放。问题二选择$\alpha=0.80$。正式策略的1月热启动固定使用0.80分位，随后连续承接库存，2—12月计入正式结果。

\subsubsection{十分钟安全反馈与费用}
令$z_t=q_t-(L_t-V_t)$表示常规电力相对实际净需求的余量。在每个交付段观察到实际需求后，采用
\begin{align}
 z_t\ge0:&\quad c_t=\min\{z_t,5000/6,(10800-S_t)/0.9\},\quad b_t=e_t=0,\quad w_t=z_t-c_t;\\
 z_t<0:&\quad b_t=\min\{-z_t,5000/6,0.9(S_t-1200)\},\quad c_t=w_t=0,\quad e_t=-z_t-b_t.
\end{align}
每步按式\eqref{eq:soc}更新库存。无外网紧急购电上限时，$e_t$消除剩余缺口，保证物理可行；它不会被用于额外给电池充电。问题二的实际账单为
\begin{equation}
 C_2=\sum_{d,t}\left(p_tg_{d,t}+5p_te_{d,t}\right).
\end{equation}
该反馈在缺口出现时尽量放电，是简单、可复现的可行控制。它没有比较跨时段紧急电价，也没有评估储存1\,kWh电能的未来价值，因此不能由可行性推导经济最优性。第\ref{sec:feedback}节另给价格感知反馈对照。
